\subsection{Sonificación y resultados}
\label{subsec:validacion_sonificacion_resultados}

La última validación consistió en comprobar la integración completa EEG--MIDI con señal real. La sesión final se identificó como \texttt{s01\_20260528} y se realizó con el montaje \texttt{ear\_eeg\_ch1\_only}, el modo \texttt{bias\_ch1\_only\_loff\_off}, \texttt{ADS\_DIAGNOSTIC\_MODE=5} y el modelo musical \texttt{modelo\_captura\_final}. El objetivo no fue demostrar una interpretación clínica de las bandas EEG, sino verificar que el sistema era capaz de adquirir señal real, calcular características, aplicar el \textit{quality gate}, transformar esas características en controles musicales y guardar eventos MIDI durante la adquisición.

La ruta validada en esta fase fue:

\begin{center}
\texttt{ADS1299 $\rightarrow$ firmware $\rightarrow$ Bridge $\rightarrow$ Python $\rightarrow$ DSP $\rightarrow$ quality gate $\rightarrow$ sonificación $\rightarrow$ notas MIDI}
\end{center}

Cada captura conservó los ficheros necesarios para reconstruir el resultado fuera de línea: \texttt{eeg\_timeseries.csv}, \texttt{metadata.json}, \texttt{windowed\_sonification\_features.csv}, \texttt{music\_snapshots.jsonl}, \texttt{music\_notes.csv} y \texttt{music\_capture\_summary.json}. Por tanto, la validación no depende únicamente de la escucha subjetiva, sino de una trazabilidad completa desde la señal EEG hasta las notas registradas.

La sesión incluyó dos precomprobaciones y cinco condiciones principales: ojos abiertos, ojos cerrados, reposo quieto, parpadeo y una repetición de ojos abiertos. Todas las capturas revisadas mantuvieron una frecuencia efectiva de \SI{250}{\hertz}, sin \textit{sample gaps} y sin estados ADS1299 inválidos. La Tabla~\ref{tab:validacion_sesion_final_musica} resume la salida musical obtenida en las condiciones principales.

\begin{table}[!htbp]
    \centering
    \small
    \caption{Resumen de la sesión final de laboratorio y salida musical registrada.}
    \label{tab:validacion_sesion_final_musica}
    \begin{tabular}{p{0.31\textwidth} c c c p{0.26\textwidth}}
        \hline
        \textbf{Condición} & \textbf{Duración} & \textbf{Diagnóstico} & \textbf{Notas} & \textbf{Lectura principal} \\
        \hline
        \texttt{eyes\_open\_rest\_60s} & \SI{60}{\second} & Dudosa & 45 & Salida musical con transitorio visible. \\
        \texttt{eyes\_closed\_rest\_60s} & \SI{60}{\second} & Dudosa & 80 & Mayor densidad de eventos, condicionada por ruido. \\
        \texttt{quiet\_rest\_60s} & \SI{60}{\second} & Dudosa & 72 & Sonificación sostenida durante reposo. \\
        \texttt{blink\_artifact\_30s} & \SI{30}{\second} & Dudosa & 34 & Respuesta ante artefacto fisiológico controlado. \\
        \texttt{eyes\_open\_repeat\_30s} & \SI{30}{\second} & Válida preliminar & 65 & Mejor evidencia de integración completa. \\
        \hline
    \end{tabular}
\end{table}

Antes de analizar cada estado, conviene aclarar cómo se interpretan las gráficas MIDI. El número de notas indica cuántos eventos quedaron registrados en \texttt{music\_notes.csv}; por tanto, una captura con más notas presenta mayor actividad musical. Si la duración de la captura es menor, la densidad de notas puede ser todavía mayor. Las notas más cortas indican eventos menos sostenidos, normalmente asociados a cambios rápidos en los controles de sonificación o a ventanas cuya calidad no se mantiene estable. El \textit{pitch} MIDI representa la altura musical: valores más altos corresponden a notas más agudas y valores más bajos a notas más graves. Por tanto, la salida musical se interpreta mediante tres aspectos principales: densidad de eventos, duración de las notas y registro musical.

A continuación se introduce cada condición antes de presentar sus láminas de resultados. Esta separación evita repetir el mismo texto dentro de cada página gráfica y permite dedicar más espacio a las figuras.

\paragraph{Ojos abiertos.}
La condición \texttt{eyes\_open\_rest\_60s} se empleó como primera captura larga con señal real. La Figura~\ref{fig:validacion_s01_01_estado} muestra la cadena completa de esta condición. La adquisición se mantiene durante \SI{60}{\second}, pero aparece un transitorio de amplitud que afecta a la interpretación de la señal; por ello, esta captura se usa como evidencia de integración real, no como EEG limpio.

\paragraph{Ojos cerrados.}
La condición \texttt{eyes\_closed\_rest\_60s} permitió observar el comportamiento del sistema en un reposo distinto. La Figura~\ref{fig:validacion_s01_02_estado} muestra una salida musical más densa que en ojos abiertos, con 80 notas registradas en \SI{60}{\second}. Esta diferencia indica que los controles musicales cambiaron, aunque el diagnóstico continúa condicionado por ruido de red y artefactos.

\paragraph{Reposo quieto.}
La condición \texttt{quiet\_rest\_60s} se planteó como estado basal de funcionamiento sostenido. La Figura~\ref{fig:validacion_s01_03_estado} muestra que el sistema mantiene adquisición, análisis, control de calidad y salida musical durante toda la captura. Se registraron 72 notas, por lo que la sonificación se mantuvo activa de forma continuada.

\paragraph{Parpadeo como artefacto controlado.}
La condición \texttt{blink\_artifact\_30s} se incluyó para comprobar la respuesta ante una señal intencionadamente contaminada. La Figura~\ref{fig:validacion_s01_04_estado} permite observar cómo el artefacto afecta a la señal temporal, a la calidad, a las bandas relativas y a los controles de sonificación. Se registraron 34 notas en \SI{30}{\second}.

\paragraph{Repetición de ojos abiertos.}
La captura \texttt{06\_eyes\_open\_repeat\_30s} fue la mejor evidencia final de integración. La Figura~\ref{fig:validacion_s01_06_estado} muestra adquisición continua, diagnóstico \texttt{valida\_preliminar}, características espectrales, control de calidad y salida MIDI clara. Se registraron 65 notas en solo \SI{30}{\second}, por lo que la densidad de eventos fue elevada respecto a las capturas de \SI{60}{\second}.

En cada lámina se mantiene el mismo orden de lectura: señal temporal, calidad de señal, bandas relativas, controles de sonificación y notas MIDI. Para ganar tamaño sin utilizar páginas apaisadas, cada estado se coloca en una página vertical propia y el bloque gráfico se ensancha ligeramente respecto al ancho normal del texto.

\clearpage
\enlargethispage{1.4cm}
\noindent\makebox[\textwidth][c]{%
\begin{minipage}{1.13\textwidth}
    \centering
    {\large\bfseries Ojos abiertos: \texttt{eyes\_open\_rest\_60s}\par}
    \vspace{0.10cm}
    \begin{minipage}[t]{0.50\linewidth}
        \centering
        \validacioninclude{width=\linewidth}{images/validacion/capturas_finales_s01_20260528/01_eyes_open_rest_60s/eeg_ch1_temporal.png}\\[-0.25em]
        {\scriptsize (a) Señal temporal CH1\par}
        \vspace{0.03cm}
        \validacioninclude{width=\linewidth}{images/validacion/capturas_finales_s01_20260528/01_eyes_open_rest_60s/calidad_senal_quality_gate.png}\\[-0.25em]
        {\scriptsize (b) Calidad de señal y \textit{quality gate}\par}
        \vspace{0.03cm}
        \validacioninclude{width=\linewidth}{images/validacion/capturas_finales_s01_20260528/01_eyes_open_rest_60s/bandpowers_relativos.png}\\[-0.25em]
        {\scriptsize (c) Bandas relativas\par}
    \end{minipage}
    \hfill
    \begin{minipage}[t]{0.48\linewidth}
        \centering
        \validacioninclude{width=\linewidth}{images/validacion/capturas_finales_s01_20260528/01_eyes_open_rest_60s/controles_sonificacion.png}\\[-0.25em]
        {\scriptsize (d) Controles de sonificación\par}
        \vspace{0.10cm}
        \validacioninclude{width=\linewidth}{images/validacion/capturas_finales_s01_20260528/01_eyes_open_rest_60s/notas_musicales.png}\\[-0.25em]
        {\scriptsize (e) Notas MIDI registradas\par}
    \end{minipage}
    \captionof{figure}{Resultado de la condición \texttt{eyes\_open\_rest\_60s}. Se muestran la señal temporal, la calidad de ventana, las bandas relativas, los controles de sonificación y las notas MIDI generadas.}
    \label{fig:validacion_s01_01_estado}
    \vspace{-0.03cm}
    {\small\noindent\textbf{Conclusión.} La captura generó 45 notas en \SI{60}{\second}. El transitorio visible confirma la necesidad de leer la música junto con el \textit{quality gate}: la salida MIDI existe y queda persistida, pero no se interpreta como EEG limpio.\par}
\end{minipage}}

\clearpage
\enlargethispage{1.4cm}
\noindent\makebox[\textwidth][c]{%
\begin{minipage}{1.13\textwidth}
    \centering
    {\large\bfseries Ojos cerrados: \texttt{eyes\_closed\_rest\_60s}\par}
    \vspace{0.10cm}
    \begin{minipage}[t]{0.50\linewidth}
        \centering
        \validacioninclude{width=\linewidth}{images/validacion/capturas_finales_s01_20260528/02_eyes_closed_rest_60s/eeg_ch1_temporal.png}\\[-0.25em]
        {\scriptsize (a) Señal temporal CH1\par}
        \vspace{0.03cm}
        \validacioninclude{width=\linewidth}{images/validacion/capturas_finales_s01_20260528/02_eyes_closed_rest_60s/calidad_senal_quality_gate.png}\\[-0.25em]
        {\scriptsize (b) Calidad de señal y \textit{quality gate}\par}
        \vspace{0.03cm}
        \validacioninclude{width=\linewidth}{images/validacion/capturas_finales_s01_20260528/02_eyes_closed_rest_60s/bandpowers_relativos.png}\\[-0.25em]
        {\scriptsize (c) Bandas relativas\par}
    \end{minipage}
    \hfill
    \begin{minipage}[t]{0.48\linewidth}
        \centering
        \validacioninclude{width=\linewidth}{images/validacion/capturas_finales_s01_20260528/02_eyes_closed_rest_60s/controles_sonificacion.png}\\[-0.25em]
        {\scriptsize (d) Controles de sonificación\par}
        \vspace{0.10cm}
        \validacioninclude{width=\linewidth}{images/validacion/capturas_finales_s01_20260528/02_eyes_closed_rest_60s/notas_musicales.png}\\[-0.25em]
        {\scriptsize (e) Notas MIDI registradas\par}
    \end{minipage}
    \captionof{figure}{Resultado de la condición \texttt{eyes\_closed\_rest\_60s}. Se muestran la señal temporal, la calidad de ventana, las bandas relativas, los controles de sonificación y las notas MIDI generadas.}
    \label{fig:validacion_s01_02_estado}
    \vspace{-0.03cm}
    {\small\noindent\textbf{Conclusión.} La captura generó 80 notas en \SI{60}{\second}, por lo que la densidad musical fue mayor que en ojos abiertos. El resultado evidencia un cambio funcional en la salida MIDI, aunque no se interpreta como prueba fisiológica concluyente por la calidad dudosa de la señal.\par}
\end{minipage}}

\clearpage
\enlargethispage{1.4cm}
\noindent\makebox[\textwidth][c]{%
\begin{minipage}{1.13\textwidth}
    \centering
    {\large\bfseries Reposo quieto: \texttt{quiet\_rest\_60s}\par}
    \vspace{0.10cm}
    \begin{minipage}[t]{0.50\linewidth}
        \centering
        \validacioninclude{width=\linewidth}{images/validacion/capturas_finales_s01_20260528/03_quiet_rest_60s/eeg_ch1_temporal.png}\\[-0.25em]
        {\scriptsize (a) Señal temporal CH1\par}
        \vspace{0.03cm}
        \validacioninclude{width=\linewidth}{images/validacion/capturas_finales_s01_20260528/03_quiet_rest_60s/calidad_senal_quality_gate.png}\\[-0.25em]
        {\scriptsize (b) Calidad de señal y \textit{quality gate}\par}
        \vspace{0.03cm}
        \validacioninclude{width=\linewidth}{images/validacion/capturas_finales_s01_20260528/03_quiet_rest_60s/bandpowers_relativos.png}\\[-0.25em]
        {\scriptsize (c) Bandas relativas\par}
    \end{minipage}
    \hfill
    \begin{minipage}[t]{0.48\linewidth}
        \centering
        \validacioninclude{width=\linewidth}{images/validacion/capturas_finales_s01_20260528/03_quiet_rest_60s/controles_sonificacion.png}\\[-0.25em]
        {\scriptsize (d) Controles de sonificación\par}
        \vspace{0.10cm}
        \validacioninclude{width=\linewidth}{images/validacion/capturas_finales_s01_20260528/03_quiet_rest_60s/notas_musicales.png}\\[-0.25em]
        {\scriptsize (e) Notas MIDI registradas\par}
    \end{minipage}
    \captionof{figure}{Resultado de la condición \texttt{quiet\_rest\_60s}. Se muestran la señal temporal, la calidad de ventana, las bandas relativas, los controles de sonificación y las notas MIDI generadas.}
    \label{fig:validacion_s01_03_estado}
    \vspace{-0.03cm}
    {\small\noindent\textbf{Conclusión.} La captura generó 72 notas en \SI{60}{\second}. El resultado confirma una sonificación sostenida durante una condición de reposo y demuestra que la cadena de adquisición, análisis, calidad y salida MIDI puede mantenerse activa durante una captura prolongada.\par}
\end{minipage}}

\clearpage
\enlargethispage{1.4cm}
\noindent\makebox[\textwidth][c]{%
\begin{minipage}{1.13\textwidth}
    \centering
    {\large\bfseries Parpadeo: \texttt{blink\_artifact\_30s}\par}
    \vspace{0.10cm}
    \begin{minipage}[t]{0.50\linewidth}
        \centering
        \validacioninclude{width=\linewidth}{images/validacion/capturas_finales_s01_20260528/04_blink_artifact_30s/eeg_ch1_temporal.png}\\[-0.25em]
        {\scriptsize (a) Señal temporal CH1\par}
        \vspace{0.03cm}
        \validacioninclude{width=\linewidth}{images/validacion/capturas_finales_s01_20260528/04_blink_artifact_30s/calidad_senal_quality_gate.png}\\[-0.25em]
        {\scriptsize (b) Calidad de señal y \textit{quality gate}\par}
        \vspace{0.03cm}
        \validacioninclude{width=\linewidth}{images/validacion/capturas_finales_s01_20260528/04_blink_artifact_30s/bandpowers_relativos.png}\\[-0.25em]
        {\scriptsize (c) Bandas relativas\par}
    \end{minipage}
    \hfill
    \begin{minipage}[t]{0.48\linewidth}
        \centering
        \validacioninclude{width=\linewidth}{images/validacion/capturas_finales_s01_20260528/04_blink_artifact_30s/controles_sonificacion.png}\\[-0.25em]
        {\scriptsize (d) Controles de sonificación\par}
        \vspace{0.10cm}
        \validacioninclude{width=\linewidth}{images/validacion/capturas_finales_s01_20260528/04_blink_artifact_30s/notas_musicales.png}\\[-0.25em]
        {\scriptsize (e) Notas MIDI registradas\par}
    \end{minipage}
    \captionof{figure}{Resultado de la condición \texttt{blink\_artifact\_30s}. Se muestran la señal temporal, la calidad de ventana, las bandas relativas, los controles de sonificación y las notas MIDI generadas.}
    \label{fig:validacion_s01_04_estado}
    \vspace{-0.03cm}
    {\small\noindent\textbf{Conclusión.} La captura generó 34 notas en \SI{30}{\second}. Su interés principal es mostrar la respuesta ante un artefacto controlado: la señal cambia, la calidad condiciona la interpretación y la música generada debe leerse como salida trazable ante una entrada contaminada.\par}
\end{minipage}}

\clearpage
\enlargethispage{1.4cm}
\noindent\makebox[\textwidth][c]{%
\begin{minipage}{1.13\textwidth}
    \centering
    {\large\bfseries Repetición de ojos abiertos: \texttt{06\_eyes\_open\_repeat\_30s}\par}
    \vspace{0.10cm}
    \begin{minipage}[t]{0.50\linewidth}
        \centering
        \validacioninclude{width=\linewidth}{images/validacion/capturas_finales_s01_20260528/06_eyes_open_repeat_30s/eeg_ch1_temporal.png}\\[-0.25em]
        {\scriptsize (a) Señal temporal CH1\par}
        \vspace{0.03cm}
        \validacioninclude{width=\linewidth}{images/validacion/capturas_finales_s01_20260528/06_eyes_open_repeat_30s/calidad_senal_quality_gate.png}\\[-0.25em]
        {\scriptsize (b) Calidad de señal y \textit{quality gate}\par}
        \vspace{0.03cm}
        \validacioninclude{width=\linewidth}{images/validacion/capturas_finales_s01_20260528/06_eyes_open_repeat_30s/bandpowers_relativos.png}\\[-0.25em]
        {\scriptsize (c) Bandas relativas\par}
    \end{minipage}
    \hfill
    \begin{minipage}[t]{0.48\linewidth}
        \centering
        \validacioninclude{width=\linewidth}{images/validacion/capturas_finales_s01_20260528/06_eyes_open_repeat_30s/controles_sonificacion.png}\\[-0.25em]
        {\scriptsize (d) Controles de sonificación\par}
        \vspace{0.10cm}
        \validacioninclude{width=\linewidth}{images/validacion/capturas_finales_s01_20260528/06_eyes_open_repeat_30s/notas_musicales.png}\\[-0.25em]
        {\scriptsize (e) Notas MIDI registradas\par}
    \end{minipage}
    \captionof{figure}{Resultado de la condición \texttt{06\_eyes\_open\_repeat\_30s}. Se muestran la señal temporal, la calidad de ventana, las bandas relativas, los controles de sonificación y las notas MIDI generadas.}
    \label{fig:validacion_s01_06_estado}
    \vspace{-0.03cm}
    {\small\noindent\textbf{Conclusión.} La captura generó 65 notas en \SI{30}{\second}, la mayor densidad de la sesión. Este resultado reúne adquisición continua, diagnóstico \texttt{valida\_preliminar}, control de calidad y salida MIDI persistida, por lo que constituye la evidencia más clara de integración EEG--MIDI.\par}
\end{minipage}}

\clearpage
La comparación entre condiciones muestra que los cambios en la señal registrada producen cambios observables en la música generada. Cuando varían la amplitud, la distribución relativa de bandas y la calidad de ventana, también cambian la densidad de notas, su duración y su registro. Esta relación no debe entenderse como una decodificación clínica del estado cerebral, sino como una validación funcional de la interfaz EEG--MIDI: el sistema transforma una señal real, con sus limitaciones experimentales, en una salida musical trazable y persistida.

En conclusión, la validación de sonificación confirma que el prototipo completa la cadena funcional desde la adquisición EEG hasta la generación de notas MIDI. La sesión \texttt{s01\_20260528} demuestra integración real, continuidad temporal, cálculo de características, aplicación del \textit{quality gate} y persistencia de la salida musical. Su valor principal es mostrar que la dinámica de la señal produce cambios observables en la música generada, manteniendo una interpretación prudente frente al ruido, los artefactos y las limitaciones del montaje experimental.
